% UNIFIED 2026-05-13
% SOURCE MAP:
%   - Plan 1/route-A-summary.tex
%       Sections 5 (closure by BDV nearly spherical), 6 (the
%       non-contradictory quantitative argument), 7 (extension to all
%       asymmetries) supply the assembly logic. We CONVERT the source
%       dimension symbol N -> n per Manuscript/00_notation_register.md.
%   - Final/NearlySphericalClosure.tex
%       Theorem 2.2 (NS), Corollary 2.4 (vol/barycenter absorption),
%       Lemmas 1.1, 1.2 (BDV Thm 3.3 and Lem 4.2(i),(iii)).
%   - Final/SaintVenantAssembly.tex
%       Theorems 2.1, 4.1 (small-alpha linear bound, sharp Saint-Venant),
%       Lemmas 1.1, 1.2, 1.3 (BDV Lem 4.2(i),(ii), Lem 5.2), and
%       Corollary 3.1 / Lemma 3.2 (conversion alpha <-> Fr).
%   - Plan 1/1306.0392v1.pdf
%       Brasco-De Philippis-Velichkov, "Faber-Krahn inequalities in sharp
%       quantitative form", arXiv:1306.0392v1 (2013).
%       Used as black boxes: Definition 4.1 (smooth asymmetry),
%       Theorem 3.3 (sharp nearly spherical), Lemma 4.2 (i),(ii),(iii),
%       Lemma 5.2 (suboptimal Fraenkel-quartic Saint-Venant),
%       Lemmas 4.9, 4.12, 4.16, Theorem 4.18 (regularity package used
%       upstream via Agents 5-6; not re-invoked here).
%   - Manuscript/00_notation_register.md
%       Sections 1-6, 14: canonical n, R, B^*, T(.), E(.), delta_T,
%       Fraenkel \mathcal A, BDV smooth alpha (renamed \alpha_BDV in
%       passages where coarea alpha(t) could collide), constants tracked
%       with explicit (n,R) dependence.
%   - Manuscript/source_map.md
%       Entries D3.1-D3.8 (nearly spherical closure and bounded sharp
%       Saint-Venant), with upstream dependencies D1.x (selection,
%       Agent 5) and D2.x (graph entry / Schauder, Agent 6).
%
% UPSTREAM HYPOTHESES (Agents 5 and 6) USED VERBATIM:
%   - (S) Quantitative selection theorem with constants
%         C_sel(n,R), q_sel(n,R), eps_sel(n,R), tau_reg(n,R).
%   - (G) Graph entry: alpha_graph(n,R), C_H''(n,R), gamma in (0,1).
%   - (I) Schauder/interpolation entry into the C^{2,gamma_0} ball:
%         alpha_sph(n,R,gamma_0), delta_sph(n,gamma_0).
%
% This subsection produces the output the next agent (Agent 8) needs:
%         bounded sharp Saint-Venant inequality
%            E(Omega)-E(B^*) >= c_SV(n,R) \mathcal A(Omega)^2
%         for Omega open, Omega \subset B_R, |Omega| = |B_1| = omega_n,
%         equivalently
%            \mathcal A(Omega)^2 <= C(n,R) delta_T(Omega).

% ---------- Subsection: Nearly-spherical closure and bounded sharp Saint--Venant ----------

\subsection{Nearly-spherical closure and bounded sharp Saint--Venant}
\label{sec:plan1-nearly-spherical}

This subsection closes the chain
\[
\text{selection}\ \longrightarrow\ \text{graph entry}\ \longrightarrow\
\text{Schauder/interpolation entry}\ \longrightarrow\
\text{BDV nearly spherical}
\]
and converts the closing inequality into the bounded sharp
Saint--Venant estimate
\begin{equation}\label{eq:plan1-target}
   \Fr(\Om)^{2}\ \le\ C(n,R)\,\deltaT(\Om),
   \qquad \Om\subset B_{R},\ |\Om|=\omega_{n},
\end{equation}
which is the input to the bounded reduction and the Kohler--Jobin
transfer of Agent~8. The argument is \emph{not} by contradiction: every
quantitative step is a substitution of upstream constants into a chain
of explicit BDV black boxes. The constants are tracked through the
section.

Throughout the subsection $n\ge 2$ and $R\ge 2$ are fixed, all sets are
bounded open subsets of $\mathbb{R}^{n}$, and we keep the notation
register conventions. (Domain-class reconciliation. The selection
principle of Agent~5 produces a quasi-open minimiser
$U_*\subset\overline{B_R}$ in the BDV sense
\cite[\S2]{BDV}; after the BDV regularity package
\cite[Lem.~4.9,~4.12,~4.16,~Thm.~4.18]{BDV} consumed by Agent~6,
$U_*$ is open up to a set of zero capacity, and all load-bearing
functionals $\deltaT,\Fr,\BDValpha,\lambda_1$ are insensitive to
capacity-zero modifications. Consequently the quasi-open class of
Agent~5 and the open class of the notation register \S2 coincide for
every quantity used below; we treat $U_*$ as an open subset of
$\overline{B_R}$ without further mention.)
The notation register conventions are: $E(\Om)=-T(\Om)/2$ is the BDV-normalised
Saint--Venant energy, $\deltaT(\Om)=E(\Om)-E(B^{*})$ is the Saint--Venant
deficit at fixed volume, $\Fr(\Om)$ is the Fraenkel asymmetry, and the
BDV smooth asymmetry is
\[
\BDValpha(\Om)\ :=\ \int_{\Om\,\triangle\, B_{1}(x_{\Om})}
   \bigl|\,1-|x-x_{\Om}|\,\bigr|\,dx,
   \qquad
   x_{\Om}=\frac{1}{|\Om|}\int_{\Om} x\,dx
\]
(\cite[Definition~4.1]{BDV}, with $n$ in place of the source
symbol $N$). The torsion quotient at the BDV smooth asymmetry is
\[
Q_{\BDValpha}(\Om)\ :=\ \frac{E(\Om)-E(B_{1})}{\BDValpha(\Om)},
   \qquad |\Om|=\omega_{n},\ \BDValpha(\Om)>0.
\]

\subsubsection{The three BDV black boxes}
\label{sec:plan1-bdv-blackboxes}

The closure uses exactly three results from
\cite{BDV}, all of which are proved there constructively (no
compactness, no selection).

\paragraph{(B1) Sharp nearly-spherical stability \textnormal{(\cite[Theorem~3.3]{BDV})}.}
For every $0<\gamma_{0}\le 1$ there exists
$\delta_{1}=\delta_{1}(n,\gamma_{0})>0$ such that the following holds.
Let $U\subset\mathbb{R}^{n}$ be open and bounded, and suppose there
exists $g\in C^{2,\gamma_{0}}(\partial B_{1})$ with
\begin{equation}\label{eq:plan1-ns-hyp}
   \partial U=\bigl\{(1+g(y))\,y:\,y\in\partial B_{1}\bigr\},
   \quad
   \|g\|_{C^{2,\gamma_{0}}(\partial B_{1})}\le\delta_{1},
   \quad
   |U|=|B_{1}|,
   \quad
   x_{U}=0.
\end{equation}
Then, with $E$ the Saint--Venant energy of
Notation~Register~\S3 and $H(g)$ the harmonic extension of $g$ to
$B_{1}$,
\begin{equation}\label{eq:plan1-BDV33}
   E(U)-E(B_{1})\ \ge\
   c_{\rm sph}(n)\,\|g\|_{H^{1/2}(\partial B_{1})}^{2},
   \qquad
   c_{\rm sph}(n):=\frac{1}{32\,n^{2}},
\end{equation}
where $\|g\|_{H^{1/2}(\partial B_{1})}^{2}=\int_{\partial B_{1}}g^{2}\,d\mathcal{H}^{n-1}+\int_{B_{1}}|\nabla H(g)|^{2}\,dx$
(\cite[Definition~3.2]{BDV}). The constant $1/(32n^{2})$ comes from
the gap to the second Steklov eigenvalue of $B_{1}$ (equal to $2$ in
every dimension), see \cite[(3.9)--(3.10)]{BDV}.

\paragraph{(B2) Smooth-asymmetry comparisons \textnormal{(\cite[Lemma~4.2 (i),(ii),(iii)]{BDV})}.}
There exist $C_{1}=C_{1}(n)>0$, $C_{2}=C_{2}(R)>0$,
$\delta_{3}=\delta_{3}(n)>0$, $C_{3}=C_{3}(n)>0$ such that:
\begin{enumerate}
\item[\textup{(i)}]
   for every bounded open $\Om\subset\mathbb{R}^{n}$ with
   $|\Om|=|B_{1}|$,
   \begin{equation}\label{eq:plan1-42i}
      C_{1}(n)\,\BDValpha(\Om)\ \ge\ |\Om\triangle B_{1}(x_{\Om})|^{2}
      \ \ge\ |B_{1}|^{2}\Fr(\Om)^{2};
   \end{equation}
\item[\textup{(ii)}]
   for every $R\ge 2$ and every bounded open
   $\Om_{1},\Om_{2}\subset B_{R}$,
   \begin{equation}\label{eq:plan1-42ii}
      \bigl|\BDValpha(\Om_{1})-\BDValpha(\Om_{2})\bigr|
      \ \le\ C_{2}(R)\,|\Om_{1}\triangle\Om_{2}|;
   \end{equation}
   in particular, choosing $\Om_{2}$ to be the Fraenkel-optimal ball for
   $\Om_{1}$ (so $\BDValpha(\Om_{2})=0$), one has, for
   $|\Om|=\omega_{n}$ and $\Om\subset B_{R}$,
   \begin{equation}\label{eq:plan1-42ii-Fr}
      \BDValpha(\Om)\ \le\ C_{2}(R)\,\omega_{n}\,\Fr(\Om);
   \end{equation}
\item[\textup{(iii)}]
   for every nearly-spherical set $U$ parametrised by $g$ with
   $\|g\|_{L^{\infty}(\partial B_{1})}\le\delta_{3}$,
   \begin{equation}\label{eq:plan1-42iii}
      \BDValpha(U)\ \le\ C_{3}(n)\,\|g\|_{L^{2}(\partial B_{1})}^{2}.
   \end{equation}
\end{enumerate}
All three estimates are proved in \cite[\S4]{BDV} by direct
rearrangement / Taylor expansion of $\BDValpha$; they are constructive.

\paragraph{(B3) Suboptimal Fraenkel-quartic Saint--Venant
\textnormal{(\cite[Lemma~5.2]{BDV})}.}
There exists $C_{8}=C_{8}(n)>0$ such that for every bounded open
$\Om\subset\mathbb{R}^{n}$ with $0<|\Om|<\infty$,
\begin{equation}\label{eq:plan1-52-raw}
   E(\Om)\,|\Om|^{-(n+2)/n}-E(B_{1})\,|B_{1}|^{-(n+2)/n}
   \ \ge\ \frac{1}{C_{8}(n)}\,\Fr(\Om)^{4}.
\end{equation}
The proof in \cite[\S5]{BDV} is by rearrangement and
P\'olya--Szeg\H{o} (no contradiction, no selection). After absorbing the
$\omega_{n}^{(n+2)/n}$ factor at fixed volume $|\Om|=\omega_{n}$,
\begin{equation}\label{eq:plan1-52-fixed}
   E(\Om)-E(B_{1})\ \ge\
   C_{8}'(n)^{-1}\,\Fr(\Om)^{4},
   \qquad C_{8}'(n):=C_{8}(n)/\omega_{n}^{(n+2)/n}.
\end{equation}

\medskip
\textbf{Important note on hypotheses.} (B1) is stated for sets that are
\emph{simultaneously} $C^{2,\gamma_{0}}$-parametrised, volume-normalised
to $|B_{1}|$, and barycentred at the origin. The three normalisations
are all load-bearing for~\eqref{eq:plan1-BDV33}: dropping any one of
them changes the form of the second-variation calculation in
\cite[\S3]{BDV}. Sections~\ref{sec:plan1-feed-graph}
and~\ref{sec:plan1-vol-bary} show that the upstream chain enforces all
three.

\subsubsection{Trace--Poincar\'e step and the pointwise
\texorpdfstring{$\alpha$}{alpha}-form}
\label{sec:plan1-alpha-form}

Combining (B1) and (B2)(iii) gives the pointwise
$Q_{\BDValpha}$-floor on nearly-spherical sets. The intermediate
$H^{1/2}\!\to L^{2}$ comparison is the trace--Poincar\'e inequality on
$\partial B_{1}$:
\begin{equation}\label{eq:plan1-poincare}
   \|g\|_{L^{2}(\partial B_{1})}^{2}\ \le\
   \|g\|_{H^{1/2}(\partial B_{1})}^{2},
\end{equation}
which is immediate from the definition of the $H^{1/2}$-norm (its $L^{2}$
summand bounds $\|g\|_{L^{2}}^{2}$ above by itself).

\begin{theorem}[Nearly-spherical closure, BDV smooth-asymmetry form]
\label{thm:plan1-NS}
Fix $0<\gamma_{0}\le 1$ and let
\begin{equation}\label{eq:plan1-deltasph}
   \delta_{\rm sph}(n,\gamma_{0})\ :=\ \min\bigl(\delta_{1}(n,\gamma_{0}),\,\delta_{3}(n)\bigr).
\end{equation}
Let $U\subset\mathbb{R}^{n}$ be open and bounded, with
$\partial U=\{(1+g(y))y:y\in\partial B_{1}\}$,
$g\in C^{2,\gamma_{0}}(\partial B_{1})$, and
\begin{equation}\label{eq:plan1-NS-hyp}
   \|g\|_{C^{2,\gamma_{0}}(\partial B_{1})}\le\delta_{\rm sph}(n,\gamma_{0}),
   \qquad |U|=|B_{1}|,
   \qquad x_{U}=0.
\end{equation}
Then
\begin{equation}\label{eq:plan1-NS-alpha}
   E(U)-E(B_{1})\ \ge\ c_{*}(n)\,\BDValpha(U),
   \qquad
   c_{*}(n)\ :=\ \frac{c_{\rm sph}(n)}{C_{3}(n)}\ =\
   \frac{1}{32\,n^{2}\,C_{3}(n)},
\end{equation}
and consequently $Q_{\BDValpha}(U)\ge c_{*}(n)$.
\end{theorem}

\begin{proof}
By~\eqref{eq:plan1-deltasph} both
hypotheses~\eqref{eq:plan1-ns-hyp} (with $\delta_{1}$) and the smallness
hypothesis of (B2)(iii) (with $\delta_{3}$) hold simultaneously, since
$\|g\|_{L^{\infty}}\le\|g\|_{C^{2,\gamma_{0}}}\le\delta_{\rm sph}\le\delta_{3}$.
Chain (B1), \eqref{eq:plan1-poincare} and (B2)(iii):
\[
   E(U)-E(B_{1})
   \overset{\eqref{eq:plan1-BDV33}}{\ge}
   c_{\rm sph}(n)\,\|g\|_{H^{1/2}}^{2}
   \overset{\eqref{eq:plan1-poincare}}{\ge}
   c_{\rm sph}(n)\,\|g\|_{L^{2}}^{2}
   \overset{\eqref{eq:plan1-42iii}}{\ge}
   \frac{c_{\rm sph}(n)}{C_{3}(n)}\,\BDValpha(U).
\]
Dividing by $\BDValpha(U)>0$ gives $Q_{\BDValpha}(U)\ge c_{*}(n)$.
\end{proof}

Theorem~\ref{thm:plan1-NS} is the pointwise (per-set) closure: no
sequence, no limit, no smallness of $Q_{\BDValpha}$ is assumed. It is
exactly the building block called $\NS$ in
\texttt{Final/NearlySphericalClosure.tex}.

\subsubsection{Feeding the upstream chain into the
nearly-spherical hypothesis}
\label{sec:plan1-feed-graph}

We now show that the upstream blocks (S), (G), (I) of Agents~5 and~6
produce, from \emph{any} $\Om_{0}\subset B_{R}$ with $|\Om_{0}|=\omega_{n}$
and $\BDValpha(\Om_{0})$ small, a set $\Otilde$ that meets all three
hypotheses of~\eqref{eq:plan1-NS-hyp}.

\paragraph{Recap of the upstream output (Agents 5, 6).}
We use the following from Agent~5 (selection) and Agent~6
(graph entry / Schauder interpolation), each stated with the constants
they fix:

\begin{itemize}
\item[(S)] (Quantitative selection, Agent~5, Theorem D1.2--D1.5 of
   \texttt{Manuscript/source\_map.md}; \texttt{Plan 1/quantitative-selection-principle.tex},
   \texttt{Final/SelectionPrinciple.tex}.)
   There exist $C_{\rm sel}=C_{\rm sel}(n,R)>0$,
   $q_{\rm sel}=q_{\rm sel}(n,R)>0$, $\eps_{\rm sel}=\eps_{\rm sel}(n,R)>0$
   such that for every open $\Om_{0}\subset B_{R}$ with
   $|\Om_{0}|=\omega_{n}$, $\BDValpha(\Om_{0})\le\eps_{\rm sel}$,
   $Q_{\BDValpha}(\Om_{0})<q_{\rm sel}$, the single-set penalised
   functional of Agent~5 has a minimiser $U_{*}\subset B_{R}$, whose
   volume-normalised companion
   \(
      \Otilde\ :=\ r_{*}^{-1}\,U_{*},
      \quad r_{*}=\bigl(|U_{*}|/|B_{1}|\bigr)^{1/n},
   \)
   satisfies $\Otilde\subset B_{R}$ (up to enlarging $R$ by a dimensional
   amount), $|\Otilde|=\omega_{n}$,
   \begin{equation}\label{eq:plan1-S-alpha}
      \tfrac{1}{2}\,\BDValpha(\Om_{0})\ \le\ \BDValpha(\Otilde)\ \le\
      \tfrac{3}{2}\,\BDValpha(\Om_{0}),
   \end{equation}
   \begin{equation}\label{eq:plan1-S-Q}
      Q_{\BDValpha}(\Otilde)\ \le\ 2\,C_{\rm sel}(n,R)\,
      Q_{\BDValpha}(\Om_{0}),
   \end{equation}
   and $\Otilde$ inherits the BDV regularity package: density estimates,
   Lipschitz nondegeneracy, and a smooth Bernoulli coefficient
   (\cite[Lemmas~4.9, 4.12, 4.16]{BDV}).

\item[(G)] (Graph entry, Agent~6,
   Theorem D2.3 of \texttt{Manuscript/source\_map.md};
   \texttt{Final/GraphEntry.tex}, \texttt{Plan 1/graph-entry-closure.tex}.)
   There exist $\alpha_{\rm graph}=\alpha_{\rm graph}(n,R)>0$ and
   $C_{H}''=C_{H}''(n,R)>0$ and $\gamma=\gamma(n,R)\in(0,1)$ such that
   every $\Otilde$ as in (S) with
   $\BDValpha(\Otilde)\le\alpha_{\rm graph}$ is, after a single
   translation $\Otilde\mapsto\Otilde-x_{\Otilde}$ by its
   barycentre, a $C^{1,\gamma}$ normal graph over $\partial B_{1}$:
   \begin{equation}\label{eq:plan1-G}
      \partial\Otilde=\bigl\{(1+g(\theta))\,\theta:\theta\in\partial B_{1}\bigr\},
      \qquad
      \|g\|_{L^{\infty}(\partial B_{1})}\le C_{H}''(n,R)\,
      \BDValpha(\Otilde)^{1/(2n)}.
   \end{equation}

\item[(I)] (Schauder interpolation, Agent~6,
   Lemma D2.7 of \texttt{Manuscript/source\_map.md};
   \texttt{Final/SchauderInterpolation.tex}.)
   There exists $\alpha_{\rm sph}=\alpha_{\rm sph}(n,R,\gamma_{0})>0$
   such that whenever (G) applies and
   $\BDValpha(\Otilde)\le\alpha_{\rm sph}$, the parametrisation $g$ of
   (G) satisfies
   \begin{equation}\label{eq:plan1-I}
      \|g\|_{C^{2,\gamma_{0}}(\partial B_{1})}\ \le\
      \delta_{\rm sph}(n,\gamma_{0}),
   \end{equation}
   the smallness threshold of Theorem~\ref{thm:plan1-NS}.
\end{itemize}

The thresholds $\eps_{\rm sel}, \alpha_{\rm graph}, \alpha_{\rm sph}$
appearing in (S), (G), (I) are independent of $\Om_{0}$ and are fixed
once $(n,R,\gamma_{0})$ is chosen. Define the small-asymmetry threshold
\begin{equation}\label{eq:plan1-asmall}
   \alpha_{\rm small}(n,R)\ :=\
   \tfrac{1}{2}\,\min\bigl\{
      \eps_{\rm sel}(n,R),\ \alpha_{\rm graph}(n,R),\ \alpha_{\rm sph}(n,R,\gamma_{0})
   \bigr\}\ >\ 0.
\end{equation}

\subsubsection{Volume and barycentre normalisation}
\label{sec:plan1-vol-bary}

The set $U_{*}$ produced by (S) is the minimiser of a penalised
functional whose volume penalty does \emph{not} force $|U_{*}|=\omega_{n}$
exactly; the energy gain forces
\begin{equation}\label{eq:plan1-vol-defect}
   \bigl||U_{*}|-\omega_{n}\bigr|\ \le\ C(n,R)\,\deltaT(U_{*}),
\end{equation}
see \cite[Lemma~4.5]{BDV} and Agent~5. We normalise both the
volume and the barycentre as follows.

\paragraph{Volume normalisation.}
The companion
\(
   \Otilde=r_{*}^{-1}U_{*},
   \quad
   r_{*}=(|U_{*}|/\omega_{n})^{1/n},
\)
satisfies $|\Otilde|=\omega_{n}$ by construction, and the torsion
equation rescales as $u_{\Otilde}(x)=r_{*}^{-2}u_{U_{*}}(r_{*}x)$, so
that $E(\Otilde)=r_{*}^{-(n+2)}E(U_{*})$ (cf.\ the rescaling computation
of \cite[\S2]{BDV}, with $N$ replaced by $n$). The Saint--Venant
deficit transforms as
\(
   \deltaT(\Otilde)=r_{*}^{-(n+2)}\bigl(E(U_{*})-r_{*}^{n+2}E(B_{1})\bigr)
\)
and~\eqref{eq:plan1-vol-defect} gives $|r_{*}-1|\le C(n,R)\,\deltaT(U_{*})$;
the energy errors are absorbed into a multiplicative
$1+O(\deltaT(U_{*}))$ factor in (S), which the constants
$C_{\rm sel}(n,R)$ already include (Agent~5).

\paragraph{Barycentre normalisation.}
After (G) provides the parametrisation $\partial\Otilde=\{(1+g)y\}$, the
barycentre $x_{\Otilde}$ and the volume normalisation $|\Otilde|=\omega_{n}$
give the two scalar identities (\cite[(3.6)--(3.7)]{BDV}, $N\to n$)
\begin{equation}\label{eq:plan1-vol-eq}
   |\Otilde|=\int_{\partial B_{1}}\frac{(1+g)^{n}}{n}\,d\mathcal{H}^{n-1}=\omega_{n},
\end{equation}
\begin{equation}\label{eq:plan1-bary-eq}
   x_{\Otilde}=\int_{\partial B_{1}}y\,\frac{(1+g)^{n+1}}{n+1}\,d\mathcal{H}^{n-1}.
\end{equation}
Expanding in $g$ and using $\|g\|_{C^{2,\gamma_{0}}}\le\delta_{\rm sph}$
yields the BDV moment bounds
(\cite[(3.6)--(3.7)]{BDV}, with $N\to n$):
\begin{equation}\label{eq:plan1-moment}
   \Bigl|\int_{\partial B_{1}}g\,d\mathcal{H}^{n-1}\Bigr|
   + \sum_{i=1}^{n}\Bigl|\int_{\partial B_{1}}y_{i}\,g\,d\mathcal{H}^{n-1}\Bigr|
   \ \le\ C(n)\,\delta_{\rm sph}\,\|g\|_{H^{1/2}(\partial B_{1})}.
\end{equation}
The right-hand side of~\eqref{eq:plan1-moment} is a controlled fraction
of $\|g\|_{H^{1/2}}^{2}$ once $\|g\|_{H^{1/2}}\ge c(n)\,\delta_{\rm sph}$;
when $\|g\|_{H^{1/2}}$ is even smaller, both sides of
\eqref{eq:plan1-BDV33} are subdominant and the small-deficit
case~\eqref{eq:plan1-NS-alpha} holds trivially after a single homothety
and translation. The detailed reduction, which is a one-step rescaling
plus translation, is recorded in
\texttt{Final/NearlySphericalClosure.tex} (Corollary~2.4) and shows
that we may replace \eqref{eq:plan1-NS-hyp} by
\begin{equation}\label{eq:plan1-NS-hyp-relaxed}
   \bigl|\,|\Otilde|-\omega_{n}\,\bigr|+|x_{\Otilde}|\ \le\
   \eta(n,\gamma_{0})\,\|g\|_{L^{2}(\partial B_{1})}
\end{equation}
at the cost of replacing $c_{*}(n)$ by
$\widetilde{c}_{*}(n):=c_{*}(n)/4$. The post-volume normalisation
$|\Otilde|=\omega_{n}$ and the barycentre translation
$\Otilde\mapsto\Otilde-x_{\Otilde}$ make this condition trivially
satisfied with $\eta=0$, so we may apply Theorem~\ref{thm:plan1-NS}
\emph{verbatim} (with $c_{*}(n)$, not $\widetilde{c}_{*}(n)$).

\medskip\noindent
\textbf{Summary of normalisations enforced.}
\begin{itemize}
   \item Volume: $|\Otilde|=\omega_{n}$ exactly, by the homothety
      $\Otilde=r_{*}^{-1}U_{*}$.
   \item Barycentre: $x_{\Otilde}=0$, by the single translation
      $\Otilde\mapsto\Otilde-x_{\Otilde}$, which preserves
      $E,\BDValpha,\Fr$ and the $C^{2,\gamma_{0}}$-norm of $g$.
   \item $C^{2,\gamma_{0}}$-smallness: $\|g\|_{C^{2,\gamma_{0}}}\le\delta_{\rm sph}$
      by (I), provided $\BDValpha(\Otilde)\le\alpha_{\rm sph}$.
\end{itemize}
None of the three normalisations is silently dropped; each is enforced
by a named operation (rescale, translate, threshold check).

\subsubsection{Small-asymmetry linear bound}
\label{sec:plan1-small-alpha}

We now carry through the implication
\((S)\to(G)\to(I)\to\text{Theorem~\ref{thm:plan1-NS}}\) explicitly, with
no compactness step. The output is the small-$\BDValpha$ linear-in-$\BDValpha$
torsion bound.

\begin{theorem}[Small-$\BDValpha$ linear Saint--Venant bound]
\label{thm:plan1-small-alpha}
Fix $0<\gamma_{0}\le 1$. Set
\begin{equation}\label{eq:plan1-qstar}
   q_{*}(n,R)\ :=\
   \frac{c_{*}(n)}{2\,C_{\rm sel}(n,R)}
   \ =\ \frac{1}{64\,n^{2}\,C_{3}(n)\,C_{\rm sel}(n,R)}.
\end{equation}
For every open $\Om_{0}\subset B_{R}$ with $|\Om_{0}|=\omega_{n}$ and
$\BDValpha(\Om_{0})\le\alpha_{\rm small}(n,R)$,
\begin{equation}\label{eq:plan1-small-alpha}
   E(\Om_{0})-E(B_{1})\ \ge\ q_{*}(n,R)\,\BDValpha(\Om_{0}).
\end{equation}
\end{theorem}

\begin{proof}
\emph{Case A:} $Q_{\BDValpha}(\Om_{0})\ge q_{\rm sel}(n,R)$. Then directly
$E(\Om_{0})-E(B_{1})\ge q_{\rm sel}\,\BDValpha(\Om_{0})\ge q_{*}\,\BDValpha(\Om_{0})$,
after shrinking $q_{\rm sel}$ if necessary so that $q_{*}\le q_{\rm sel}$.

\emph{Case B:} $Q_{\BDValpha}(\Om_{0})<q_{\rm sel}(n,R)$. Since
$\BDValpha(\Om_{0})\le\alpha_{\rm small}\le\eps_{\rm sel}$, the
selection (S) applies and yields $\Otilde$ satisfying
\eqref{eq:plan1-S-alpha}--\eqref{eq:plan1-S-Q}. In particular
\[
   \BDValpha(\Otilde)\ \le\
   \tfrac{3}{2}\,\BDValpha(\Om_{0})
   \ \le\ \tfrac{3}{2}\,\alpha_{\rm small}(n,R)
   \ \le\ \min\{\alpha_{\rm graph}(n,R),\,\alpha_{\rm sph}(n,R,\gamma_{0})\},
\]
where the last inequality holds by~\eqref{eq:plan1-asmall} (the factor
$\tfrac{1}{2}$ in $\alpha_{\rm small}$ absorbs the factor $\tfrac{3}{2}$
from (S2)). Thus (G) gives the $C^{1,\gamma}$-graph
\eqref{eq:plan1-G}, and (I) gives the $C^{2,\gamma_{0}}$-smallness
\eqref{eq:plan1-I}.

After the single barycentre translation
$\Otilde\mapsto\Otilde-x_{\Otilde}$ (which preserves all relevant
quantities, see \S\ref{sec:plan1-vol-bary}), the
hypotheses~\eqref{eq:plan1-NS-hyp} of Theorem~\ref{thm:plan1-NS} are
satisfied. Therefore
\[
   Q_{\BDValpha}(\Otilde)\ \ge\ c_{*}(n).
\]
Combined with~\eqref{eq:plan1-S-Q},
\[
   c_{*}(n)\ \le\ Q_{\BDValpha}(\Otilde)\ \le\
   2\,C_{\rm sel}(n,R)\,Q_{\BDValpha}(\Om_{0}),
\]
whence
\(
   Q_{\BDValpha}(\Om_{0})\ \ge\ c_{*}(n)/(2C_{\rm sel}(n,R))=q_{*}(n,R),
\)
which rearranges to~\eqref{eq:plan1-small-alpha}.
\end{proof}

\begin{remark}[Where the BDV contradiction step appears, and what
replaces it]\label{rem:plan1-no-contradiction}
At no point of the proof of Theorem~\ref{thm:plan1-small-alpha} did we
assume $Q_{\BDValpha}(\Om_{0})$ is arbitrarily small. The two cases
(``$Q_{\BDValpha}\ge q_{\rm sel}$'' and ``$Q_{\BDValpha}<q_{\rm sel}$'')
are an explicit dichotomy with thresholds depending on $(n,R)$. In the
second case the conclusion comes from a single application of (S), (G),
(I), and Theorem~\ref{thm:plan1-NS}. The original BDV proof
(\cite[Theorem~4.3]{BDV}) is by contradiction with a minimising
sequence; here it is replaced by the single-set selection of Agent~5.
The conversion of the contradiction step required no additional
expansion beyond Agent~5's quotient-preservation bound
\eqref{eq:plan1-S-Q}: the chain (S)+(G)+(I)+NS contracts the
contradiction into one arithmetic comparison
$Q_{\BDValpha}(\Otilde)\le 2C_{\rm sel}Q_{\BDValpha}(\Om_{0})$.
\end{remark}

\subsubsection{Conversion to the Fraenkel asymmetry}
\label{sec:plan1-conversion}

We now convert~\eqref{eq:plan1-small-alpha} into a Fraenkel-squared
lower bound on $\deltaT$. This is the step that needs both directions
of BDV Lemma~4.2 ((B2)(i) and (B2)(ii)): the linear comparison
\eqref{eq:plan1-42i} produces the sharp exponent~$2$, and the Lipschitz
comparison \eqref{eq:plan1-42ii-Fr} ensures that the small-$\BDValpha$
hypothesis can be entered from a smallness assumption on $\Fr$ itself.

\begin{corollary}[Small-$\Fr$ quadratic torsion bound]
\label{cor:plan1-small-Fr}
Define
\begin{equation}\label{eq:plan1-anear}
   a_{\rm near}(n,R)\ :=\
   \frac{\alpha_{\rm small}(n,R)}{C_{2}(R)\,\omega_{n}}\ >\ 0,
   \qquad
   \kappa_{\rm near}(n,R)\ :=\
   \frac{q_{*}(n,R)\,\omega_{n}^{2}}{C_{1}(n)}.
\end{equation}
For every open $\Om\subset B_{R}$ with $|\Om|=\omega_{n}$ and
$\Fr(\Om)\le a_{\rm near}(n,R)$,
\begin{equation}\label{eq:plan1-small-Fr}
   E(\Om)-E(B_{1})\ \ge\ \kappa_{\rm near}(n,R)\,\Fr(\Om)^{2}.
\end{equation}
\end{corollary}

\begin{proof}
By~\eqref{eq:plan1-42ii-Fr},
$\BDValpha(\Om)\le C_{2}(R)\,\omega_{n}\,\Fr(\Om)\le C_{2}\,\omega_{n}\,a_{\rm near}=\alpha_{\rm small}(n,R)$.
Theorem~\ref{thm:plan1-small-alpha} therefore gives
$E(\Om)-E(B_{1})\ge q_{*}(n,R)\,\BDValpha(\Om)$. By~\eqref{eq:plan1-42i},
$\BDValpha(\Om)\ge|B_{1}|^{2}\Fr(\Om)^{2}/C_{1}(n)=\omega_{n}^{2}\Fr(\Om)^{2}/C_{1}(n)$.
Combining,
$E(\Om)-E(B_{1})\ge (q_{*}\omega_{n}^{2}/C_{1})\,\Fr(\Om)^{2}=\kappa_{\rm near}\,\Fr(\Om)^{2}$.
\end{proof}

\subsubsection{Far-asymmetry regime}
\label{sec:plan1-far}

The far regime $\Fr(\Om)\ge a_{\rm near}$ is handled by the suboptimal
$\Fr^{4}$ bound (B3); no contradiction is needed in this regime either.

\begin{lemma}[Far-$\Fr$ quadratic torsion bound]\label{lem:plan1-far}
For every $a>0$ and every open $\Om\subset B_{R}$ with
$|\Om|=\omega_{n}$ and $\Fr(\Om)\ge a$,
\begin{equation}\label{eq:plan1-far}
   E(\Om)-E(B_{1})\ \ge\ C_{8}'(n)^{-1}\,\Fr(\Om)^{4}
   \ \ge\ \frac{a^{2}}{C_{8}'(n)}\,\Fr(\Om)^{2}.
\end{equation}
\end{lemma}

\begin{proof}
The first inequality is~\eqref{eq:plan1-52-fixed}; the second uses
$\Fr(\Om)^{4}=\Fr(\Om)^{2}\,\Fr(\Om)^{2}\ge a^{2}\,\Fr(\Om)^{2}$ once
$\Fr(\Om)\ge a$.
\end{proof}

\subsubsection{Bounded sharp Saint--Venant inequality}
\label{sec:plan1-SV}

\begin{theorem}[Bounded sharp Saint--Venant, Plan~1]
\label{thm:plan1-SV}
Let $n\ge 2$, $R\ge 2$, $0<\gamma_{0}\le 1$. Define
\begin{equation}\label{eq:plan1-cSV}
   c_{\rm SV}(n,R)\ :=\
   \min\!\Bigl\{
      \kappa_{\rm near}(n,R),\
      \frac{a_{\rm near}(n,R)^{2}}{C_{8}'(n)}
   \Bigr\}\ >\ 0,
\end{equation}
where $\kappa_{\rm near}, a_{\rm near}$ are given by
\eqref{eq:plan1-anear}. For every open $\Om\subset B_{R}$ with
$|\Om|=\omega_{n}=|B^{*}|$,
\begin{equation}\label{eq:plan1-SV}
   E(\Om)-E(B^{*})\ \ge\ c_{\rm SV}(n,R)\,\Fr(\Om)^{2},
\end{equation}
equivalently,
\begin{equation}\label{eq:plan1-SV-deficit}
   \Fr(\Om)^{2}\ \le\ \frac{1}{c_{\rm SV}(n,R)}\,\deltaT(\Om).
\end{equation}
The exponent $2$ on $\Fr$ is sharp: it is saturated asymptotically by
small volume-preserving ellipsoidal perturbations of $B_{1}$
(\cite[Introduction, p.~2]{BDV}).
\end{theorem}

\begin{proof}
The Fraenkel asymmetry is contained in $[0,2)$; we split it at
$a_{\rm near}(n,R)$.

\emph{Case 1 (small Fraenkel):}
$\Fr(\Om)\le a_{\rm near}(n,R)$. Corollary~\ref{cor:plan1-small-Fr}
gives $E(\Om)-E(B_{1})\ge\kappa_{\rm near}(n,R)\,\Fr(\Om)^{2}\ge c_{\rm SV}(n,R)\,\Fr(\Om)^{2}$.

\emph{Case 2 (large Fraenkel):}
$\Fr(\Om)\ge a_{\rm near}(n,R)$. Lemma~\ref{lem:plan1-far} with
$a=a_{\rm near}$ gives
$E(\Om)-E(B_{1})\ge a_{\rm near}^{2}/C_{8}'\,\Fr(\Om)^{2}\ge c_{\rm SV}(n,R)\,\Fr(\Om)^{2}$.

Since $\{\Fr(\Om)\le a_{\rm near}\}\cup\{\Fr(\Om)\ge a_{\rm near}\}=[0,2)$,
the two cases cover all admissible $\Om$, and the bound is uniform with
constant $c_{\rm SV}(n,R)$ as in~\eqref{eq:plan1-cSV}. The deficit form
\eqref{eq:plan1-SV-deficit} follows since $E(B_{1})=E(B^{*})$ when
$|\Om|=|B^{*}|=\omega_{n}$, and $\deltaT(\Om)=E(\Om)-E(B^{*})$ by the
notation register.
\end{proof}

\subsubsection{Constant chain: explicit dependences}
\label{sec:plan1-constants}

The final constant $c_{\rm SV}(n,R)$ in
Theorem~\ref{thm:plan1-SV} is the minimum of a chain of named
quantities. We record their dependences explicitly.

\renewcommand{\arraystretch}{1.15}
\begin{center}
\begin{tabular}{|l|l|l|}
\hline
Constant & Dependence & Source \\
\hline
$C_{1}(n)$ & $n$ & BDV Lem.\ 4.2(i), \eqref{eq:plan1-42i}\\
$C_{2}(R)$ & $R$ & BDV Lem.\ 4.2(ii), \eqref{eq:plan1-42ii}\\
$C_{3}(n)$ & $n$ & BDV Lem.\ 4.2(iii), \eqref{eq:plan1-42iii}\\
$\delta_{1}(n,\gamma_{0})$ & $n,\gamma_{0}$ & BDV Thm.\ 3.3, \eqref{eq:plan1-BDV33}\\
$\delta_{3}(n)$ & $n$ & BDV Lem.\ 4.2(iii)\\
$C_{8}'(n)$ & $n$ & BDV Lem.\ 5.2, \eqref{eq:plan1-52-fixed}\\
$c_{\rm sph}(n)=1/(32n^{2})$ & $n$ & \eqref{eq:plan1-BDV33}\\
$\delta_{\rm sph}(n,\gamma_{0})=\min(\delta_{1},\delta_{3})$ & $n,\gamma_{0}$ & \eqref{eq:plan1-deltasph}\\
$c_{*}(n)=c_{\rm sph}/C_{3}=1/(32n^{2}C_{3}(n))$ & $n$ & Thm.\ \ref{thm:plan1-NS}\\
$C_{\rm sel}(n,R), q_{\rm sel}(n,R), \eps_{\rm sel}(n,R)$ & $n,R$ & Agent~5 (S)\\
$\alpha_{\rm graph}(n,R), C_{H}''(n,R)$ & $n,R$ & Agent~6 (G)\\
$\alpha_{\rm sph}(n,R,\gamma_{0})$ & $n,R,\gamma_{0}$ & Agent~6 (I)\\
\hline
$\alpha_{\rm small}(n,R)$ & $n,R$ & \eqref{eq:plan1-asmall}\\
$q_{*}(n,R)=c_{*}/(2C_{\rm sel})$ & $n,R$ & \eqref{eq:plan1-qstar}\\
$\kappa_{\rm near}(n,R)=q_{*}\omega_{n}^{2}/C_{1}$ & $n,R$ & \eqref{eq:plan1-anear}\\
$a_{\rm near}(n,R)=\alpha_{\rm small}/(C_{2}\omega_{n})$ & $n,R$ & \eqref{eq:plan1-anear}\\
$c_{\rm SV}(n,R)=\min\{\kappa_{\rm near},a_{\rm near}^{2}/C_{8}'\}$ & $n,R$ & \eqref{eq:plan1-cSV}\\
\hline
\end{tabular}
\end{center}

\textbf{Pure $(n,R)$-dependence of $c_{\rm SV}$.}
Every quantity in the chain depends only on $(n,R)$, possibly through
the auxiliary parameter $\gamma_{0}\in(0,1]$ which is fixed once and for
all (we may take $\gamma_{0}=1/2$). The dependence on $R$ enters at
exactly three points:
\begin{enumerate}
\item Selection $C_{\rm sel}(n,R), q_{\rm sel}(n,R), \eps_{\rm sel}(n,R)$
   (Agent~5).
\item Graph entry $\alpha_{\rm graph}(n,R), C_{H}''(n,R)$,
   Schauder interpolation $\alpha_{\rm sph}(n,R,\gamma_{0})$
   (Agent~6).
\item Lipschitz comparison $C_{2}(R)$ (BDV Lemma 4.2(ii)).
\end{enumerate}
The nearly-spherical floor $c_{*}(n)$, the BDV
constants $C_{1}(n),C_{3}(n),C_{8}'(n)$, and the multipliers $\omega_{n}$
are purely dimensional. The constant $c_{\rm sph}(n)=1/(32n^{2})$
deteriorates polynomially in $n$.

\subsubsection{Hypothesis boundaries and degeneration of the
nearly-spherical hypothesis}
\label{sec:plan1-boundaries}

The hypotheses of the BDV nearly-spherical theorem (B1) fail in three
regimes; we record how each is handled here, without silently dropping
any hypothesis.

\paragraph{(a) $\BDValpha(\Om_{0})>\alpha_{\rm small}(n,R)$.}
The selection (S) is still applicable only down to
$\BDValpha(\Om_{0})\le\eps_{\rm sel}(n,R)$. For
$\BDValpha(\Om_{0})>\alpha_{\rm small}$ but $\Fr(\Om)$ remains in
$[0,2)$, we use the far-Fraenkel regime
(Lemma~\ref{lem:plan1-far}). Translating
$\BDValpha\ge\alpha_{\rm small}$ to a Fraenkel lower bound by
\eqref{eq:plan1-42ii-Fr} gives
$\Fr(\Om)\ge\BDValpha/(C_{2}\omega_{n})\ge a_{\rm near}$, so Case~2 of
Theorem~\ref{thm:plan1-SV} applies. Thus no $\BDValpha$-large regime is
left uncovered.

\paragraph{(b) Graph regime not reached because of (G) failure.}
The graph-entry theorem (G) fails if $\BDValpha(\Otilde)>\alpha_{\rm graph}$.
Since (S) preserves $\BDValpha$ up to a factor $\tfrac{3}{2}$ and we
chose $\alpha_{\rm small}\le\alpha_{\rm graph}/2$ in
\eqref{eq:plan1-asmall} (the factor $\tfrac{1}{2}$), Case~B of the proof
of Theorem~\ref{thm:plan1-small-alpha} never enters the regime where
(G) fails. If $\BDValpha(\Om_{0})>\alpha_{\rm small}$ we fall into the
previous paragraph (a).

\paragraph{(c) Schauder/interpolation regime not reached because of (I)
failure.}
Analogously, (I) fails only if $\BDValpha(\Otilde)>\alpha_{\rm sph}$. By
the same choice of $\alpha_{\rm small}\le\alpha_{\rm sph}/2$ we never
enter that regime in Case~B. If $\BDValpha(\Om_{0})>\alpha_{\rm small}$
we again fall into (a). Equivalently: the small-asymmetry threshold
\eqref{eq:plan1-asmall} is chosen as the smallest of the three upstream
thresholds (divided by $2$ to absorb (S2)), so all three
\emph{simultaneously} hold for Case~B.

\paragraph{(d) Sharp boundary $\BDValpha\sim\alpha_{\rm small}$.}
The proof of Theorem~\ref{thm:plan1-SV} glues the two regimes at
$\Fr=a_{\rm near}$, which by \eqref{eq:plan1-anear} corresponds to
$\BDValpha=\alpha_{\rm small}$ (via the linear inequality \eqref{eq:plan1-42ii-Fr}).
At the glue point both
$\kappa_{\rm near}(n,R)\Fr^{2}$ and $(a_{\rm near}^{2}/C_{8}')\Fr^{2}$
are positive, so $c_{\rm SV}(n,R)>0$ is the minimum of two strictly
positive quantities, with explicit lower bounds
\[
   c_{\rm SV}(n,R)\ \ge\
   \min\!\Bigl\{
      \frac{\omega_{n}^{2}}{C_{1}(n)}\cdot\frac{c_{*}(n)}{2C_{\rm sel}(n,R)},\
      \frac{\alpha_{\rm small}(n,R)^{2}}{C_{2}(R)^{2}\,\omega_{n}^{2}\,C_{8}'(n)}
   \Bigr\}.
\]
No additional smallness hypothesis is needed.

\subsubsection{Summary: what this subsection delivers}
\label{sec:plan1-summary}

Combining the above with the notation $\deltaT(\Om)=E(\Om)-E(B^{*})\ge 0$
of the register, we obtain the bounded sharp Saint--Venant inequality
\begin{equation}\label{eq:plan1-final}
   \boxed{\quad
   \Fr(\Om)^{2}\ \le\ C(n,R)\,\deltaT(\Om),
   \qquad
   C(n,R)=\frac{1}{c_{\rm SV}(n,R)},
   \quad}
\end{equation}
valid for every open $\Om\subset B_{R}$ with $|\Om|=\omega_{n}=|B^{*}|$.
The constant $C(n,R)=1/c_{\rm SV}(n,R)$ has explicit dependence on
$n,R$ traced through Theorem~\ref{thm:plan1-SV} and the constant chain
of \S\ref{sec:plan1-constants}; no compactness or contradiction
argument enters. The bounded reduction
$c_{\rm SV}(n,R_{*}(n))\Rightarrow c_{\rm SV}^{\rm glob}(n)$ that removes
the $R$-dependence is the subject of Agent~8.

\paragraph{Status statement.}
\begin{itemize}
\item Theorem~\ref{thm:plan1-NS}: assembled in this subsection from
   (B1) and (B2)(iii).
\item Theorem~\ref{thm:plan1-small-alpha}: assembled in this subsection
   from (S)+(G)+(I)+Theorem~\ref{thm:plan1-NS}.
\item Corollary~\ref{cor:plan1-small-Fr},
   Lemma~\ref{lem:plan1-far}, Theorem~\ref{thm:plan1-SV}: assembled in
   this subsection from (S)+(G)+(I)+NS, (B2)(i)+(ii), and (B3).
\item Upstream black boxes (B1)+(B2)+(B3): cited verbatim from
   \cite{BDV} as \cite[Thm.\ 3.3, Lem.\ 4.2 (i),(ii),(iii),
   Lem.\ 5.2]{BDV}.
\item Inputs (S), (G), (I) from Agents~5, 6: cited by section number
   inside \texttt{Manuscript/02\_plan1\_draft.tex} once Agent~19
   integrates the three Plan 1 subsections.
\end{itemize}

\paragraph{Contradiction step.}
The BDV proof of the small-$\BDValpha$ linear bound is by
contradiction (\cite[Theorem 4.3]{BDV}); here it is replaced by the
single-set selection of Agent~5, which compresses the
sequence/compactness/limit argument into the explicit
quotient-preservation bound (S3). No additional expansion of the source
contradiction step was needed beyond Agent~5's selection theorem:
Theorem~\ref{thm:plan1-small-alpha} is a direct combinatorial
combination of (S), (G), (I), and the pointwise nearly-spherical
inequality, performed in \S\ref{sec:plan1-small-alpha}.
